% grace_arira_2026.tex
% AI-Assisted Development of Materials for Enterprise Risk Management
% Martin F. Grace --- ARIA 2026
% All slide content drawn from aria_paper.tex (revised version)
% =======================================================================
\documentclass[aspectratio=169]{beamer}
\usepackage{tikz}
\usepackage{booktabs}
\usepackage{erm_beamer_theme}

\title{AI-Assisted Development of\\Course Materials for Enterprise Risk Management}
\subtitle{A Field Report}
\author{Martin F.\ Grace}
\institute{%
  Hanson Family Chair \& Professor of Finance\\
  Faculty Director, Vaughan Institute of Risk Management \& Insurance\\
  Tippie College of Business, University of Iowa}
\date{ARIA --- August 2026}

\begin{document}

% -----------------------------------------------------------------------
\begin{frame}[plain]
  \titlepage
\end{frame}

% -----------------------------------------------------------------------
\section{Overview}
% -----------------------------------------------------------------------

\begin{frame}{Three Guiding Ideas}

\begin{itemize}\setlength\itemsep{1em}
  \item The work was \textbf{genuinely collaborative}: the LLMs performed
        substantial drafting and formatting, but they did not set the
        intellectual agenda.
  \item The bottleneck shifted from \textbf{producing a chapter} to
        \textbf{verifying the text and writing style}---with implications
        for how we think about authors' responsibilities.
  \item The project produced \textbf{infrastructure} (style guides, R figure
        themes, a GitHub distribution) that will outlast any single LLM
        version.
\end{itemize}

\bigskip
\textit{Whether that infrastructure was worth the time is left to the
reader, and to my chair or dean.}

\end{frame}

% -----------------------------------------------------------------------
\begin{frame}{Road Map}

\begin{enumerate}\setlength\itemsep{0.8em}
  \item Why this textbook?
  \item The collaboration architecture
  \item Building a consistent style system
  \item Chapter-by-chapter production
  \item From chapters to book
  \item Quality control
  \item Implications for academic authorship
\end{enumerate}

\end{frame}

% -----------------------------------------------------------------------
\section{Why This Textbook?}
% -----------------------------------------------------------------------

\begin{frame}{The Market Gap}

Traditional undergraduate risk management texts suffer from a particular
defect stemming from the relatively small size of our market.

\bigskip
To sell at scale, a book has to be usable in multiple classes. That means
authors tend to put:

\begin{itemize}\setlength\itemsep{0.5em}
  \item chapters on life and health insurance,
  \item detailed chapters on auto and commercial property,
  \item a few chapters on risk management, and
  \item maybe one chapter that mentions ERM almost in passing.
\end{itemize}

\bigskip
\textit{ERM} has nine chapters and they are all \textbf{integrated} into
the ERM framework.

\end{frame}

% -----------------------------------------------------------------------
\begin{frame}{What the Second Edition Needed}

The first edition was usable. It was also, in places,
\begin{itemize}\setlength\itemsep{0.5em}
  \item structurally repetitive,
  \item visually sparse, and
  \item more enthusiastic about lists than about exposition.
\end{itemize}

\bigskip
Two gaps needed filling:
\begin{itemize}\setlength\itemsep{0.5em}
  \item \textbf{Financial risk management with derivatives}---previously
        implied rather than taught.
  \item \textbf{AI as an enterprise-wide risk category}---not a technology
        sidebar.
\end{itemize}

\bigskip
Nine fully revised chapters, cross-references tightened so that appetite,
portfolio correlation, financing, and governance read as \emph{one
conversation}.

\end{frame}

% -----------------------------------------------------------------------
\section{Collaboration Architecture}
% -----------------------------------------------------------------------

\begin{frame}{The Software Stack}

\small
\begin{tabular}{@{}lp{0.62\linewidth}@{}}
\toprule
\textbf{Category} & \textbf{Software} \\
\midrule
\textbf{AI authoring} &
  Claude 4.6 Sonnet (drafting, conversion, math); Claude 5 Fable
  (prose and pedagogy); Perplexity (case research);
  Cursor Composer (assembly, slides, plotting) \\[0.4em]
\textbf{Environment} &
  Cursor (project editor); Git / GitHub; Windows 10 \\[0.4em]
\textbf{Conversion} &
  Python (Word/RTF $\to$ Markdown);
  Pandoc (Markdown $\to$ \LaTeX{} / PowerPoint);
  MS Excel (figure data) \\[0.4em]
\textbf{Typesetting} &
  MiKTeX 26.x; \LaTeX{} to PDF and Beamer slides \\[0.4em]
\textbf{Figures} & R 4.6 \\
\bottomrule
\end{tabular}

\bigskip
\small
Total AI expenditure: approximately \$70--\$90 for $\sim$100 hours of
author time.

\end{frame}

% -----------------------------------------------------------------------
\begin{frame}{The Workflow --- Per Chapter}

\begin{enumerate}\setlength\itemsep{0.7em}
  \item Author supplied prior manuscript chapter, an introduction case draft
        (Perplexity as search engine), and pointers to the style guide.
  \item Model proposed questions \emph{before} large rewrites.
  \item Model produced a revised Markdown chapter \texttt{*.md}, then a
        matching \LaTeX{} \texttt{.tex} file.
  \item Author compiled with \texttt{pdflatex}, inspected the PDF, and sent
        targeted correction requests (math, tone, structure).
  \item Model updated the style guide to capture anything new.
\end{enumerate}

\end{frame}

% -----------------------------------------------------------------------
\begin{frame}{Typical Prompts}

\textbf{For a chapter opening case:}
\begin{quote}\small
``Can you suggest a real-world company example that illustrates how a
company might change its risk portfolio? I want a real example, not a
made-up one. It should be a U.S.\ company.\ Please make it no longer than
a page and a half.''
\end{quote}

\medskip
\textbf{For a chapter revision:}
\begin{quote}\small
``Can you suggest any box items that might take up half a page for class
discussion? I then want a \LaTeX{} file after the new Markdown file is
created. Before you do any of this, please think about the task and ask me
any questions you may have.''
\end{quote}

\end{frame}

% -----------------------------------------------------------------------
\begin{frame}{Division of Labor}

\begin{center}
\small
\begin{tabular}{@{}p{0.44\linewidth}p{0.44\linewidth}@{}}
\toprule
\textbf{AI assistant} & \textbf{Author} \\
\midrule
First-draft exposition and section organization
  & Outline, learning objectives, emphasis \\[0.5em]
Converting Markdown to consistent \LaTeX{}
  & Economic reasoning and regulatory nuance \\[0.5em]
Applying style environments (boxes, colors)
  & Choosing real cases and verifying facts \\[0.5em]
Generating Beamer slide skeletons
  & Lecture pacing and institution-specific framing \\[0.5em]
Flagging structural repetition
  & Deciding what to cut \\[0.5em]
Patiently regenerating after the author changes his mind
  & Explaining why the third COSO rewrite still reads like a memo \\
\bottomrule
\end{tabular}
\end{center}

\bigskip
\small
Models excel at \emph{structure} and \emph{fluency};
humans retain \emph{judgment} and \emph{accountability}.

\end{frame}

% -----------------------------------------------------------------------
\section{Building a Style System}
% -----------------------------------------------------------------------

\begin{frame}{The Style Infrastructure}

A textbook that looks like nine different books teaches students that ERM
integration is optional.

\bigskip
Deliverables accumulated chapter by chapter:

\begin{itemize}\setlength\itemsep{0.6em}
  \item \texttt{erm\_textbook\_style.sty} --- colors, box environments,
        running heads, chapter opening layout
  \item \texttt{erm\_textbook\_style\_guide.md/.tex} --- parallel human- and
        machine-readable specifications
  \item \texttt{erm\_textbook\_style.r} --- ggplot theme matching the print
        palette
  \item \texttt{intro case style guide v2.md} --- narrative rules for opening
        cases
  \item \texttt{textbook\_style\_manual.md} --- voice rules, including a
        prohibition on one-sentence paragraphs
\end{itemize}

\end{frame}

% -----------------------------------------------------------------------
\begin{frame}{Style Guides as Prompt Engineering}

\begin{block}{The key insight}
Instead of re-explaining box conventions in every session, writing
``follow \texttt{erm\_textbook\_style\_guide.md}'' produced structurally
consistent output.
\end{block}

\bigskip
When the model drifted---for example, labeling Chapter~2 as Chapter~1---the
fix was usually a one-line correction, documented so it would not recur.

\bigskip
\textit{This infrastructure is the part of the project most likely to
survive the next model upgrade. The prose will age, but the \texttt{.sty}
files will still be useful.}

\end{frame}

% -----------------------------------------------------------------------
\section{Chapter-by-Chapter Production}
% -----------------------------------------------------------------------

\begin{frame}{Chapter Highlights}

\begin{description}\setlength\itemsep{0.5em}
  \item[Ch.~2 (Frameworks)] COSO boxed, ISO~31000 in comparative prose,
        GM ignition narrative up front; style guide born.
  \item[Ch.~3 (Identification)] Enterprise-vs.-silo distinction explicit
        before the four-quadrant taxonomy; 10-K risk-factor discussion added.
  \item[Ch.~4 \& 5 (Measurement, Appetite)] LTCM/\textit{Moneyball} boxes;
        Adidas risk-appetite case; the pedagogical tension clarified.
  \item[Ch.~7 (Loss Control)] Technical appendix merged in; the
        ``crazy rounding'' incident---numbers corrected in both Markdown
        and \LaTeX{}.
  \item[Ch.~8 (Risk Finance)] Three legacy sub-chapters merged into one.
  \item[Ch.~9 (Governance)] Wells Fargo, Boeing 737 MAX, and financial
        crisis-era failures integrated.
\end{description}

\end{frame}

% -----------------------------------------------------------------------
\begin{frame}{Opening Cases: The Hardest Editorial Problem}

Each chapter opens with a narrative case in a blue box.

\bigskip
Getting the AI to write a \emph{story} rather than a memo required
persistent instruction. Early drafts ended with:

\begin{quote}
``This case illustrates the importance of enterprise risk management
frameworks, which will be explored in greater depth throughout this
chapter.''
\end{quote}

Technically true. Boring.

\bigskip
I removed those sentences aggressively. The intro-case style guide exists
because telling the model ``be engaging'' is not operational.

\end{frame}

% -----------------------------------------------------------------------
\begin{frame}{What ``More Engaging'' Actually Meant}

\begin{itemize}\setlength\itemsep{0.8em}
  \item \textbf{Narrative before taxonomy.}
        Open with a firm in trouble or tension, not a definition list.
  \item \textbf{Boxes for dense material.}
        Frameworks (COSO, ISO), regulatory notes, and pitfalls break up
        long exposition.
  \item \textbf{Cross-chapter callbacks.}
        Appetite in Chapter~5 should echo identification in Chapter~3 and
        portfolio thinking in Chapter~6.
  \item \textbf{Fewer throat-clearing transitions.}
        ``This will be explored in greater depth'' was banned.
  \item \textbf{Concrete numbers with caveats.}
        Stylized examples stay, with warnings not to treat them as literal
        policy advice.
\end{itemize}

\end{frame}

% -----------------------------------------------------------------------
\section{From Chapters to Book}
% -----------------------------------------------------------------------

\begin{frame}{Integration}

The master file \texttt{erm\_book.tex} uses \texttt{docmute} to include
standalone chapter preambles without duplication, and adds:

\begin{itemize}\setlength\itemsep{0.6em}
  \item Cover art, title and copyright pages
  \item Preface (written by the author, with honest AI disclosure)
  \item Table of contents
  \item Consolidated bibliography
  \item Index (\texttt{makeindex})
\end{itemize}

\bigskip
Claude 5 Fable handled this phase. Some human troubleshooting was required.
At one point I asked whether the agent was stuck; it was, in fact, still
running in circles.

\bigskip
\textit{Patience is a human virtue, which I do not have enough of,
especially when it comes to software.}

\end{frame}

% -----------------------------------------------------------------------
\begin{frame}{Lecture Slides and Figures}

\textbf{Beamer vs.\ PowerPoint}

\smallskip
Beamer writes slides as a structured document compiled into PDF---especially
well-suited when the chapter includes equations, technical notation, or
carefully organized prose.

\bigskip
Nine chapter slide decks were generated in parallel using a custom
\texttt{erm\_beamer\_theme.sty} matching the textbook palette: consistent
fonts, gold accents, and institute lines for the University of Iowa.

\bigskip
\textbf{Figures}

\smallskip
R scripts in \texttt{rscripts/} (e.g., J\&J/Tylenol market-share plot,
aggregate loss distributions, financing spectrum). The R theme aligns
ggplot output with \LaTeX{} colors so charts do not look like they wandered
in from a different world.

\end{frame}

% -----------------------------------------------------------------------
\section{Quality Control}
% -----------------------------------------------------------------------

\begin{frame}{Where the Human Stayed in the Loop}

\begin{block}{Facts and citations}
The AI gave a strong, concise explanation of insurance expense ratios and
included a citation: journal, volume, pages, authors---everything except
reality. The fix is verification, not prompt prettiness.

Introductory cases were constrained to verifiable public sources: 10-Ks,
securities filings, major press, or business sources.
\end{block}

\bigskip
\begin{block}{Mathematics}
Chapter~7's appendix required checking inequalities, present-value
expressions, and worked examples. Students expect to be able to work it
out and follow the math. AI does not think that is so important. This
creates emails: ``Something doesn't add up! The book seems wrong.''
\end{block}

\end{frame}

% -----------------------------------------------------------------------
\begin{frame}{Voice and Pedagogy}

\begin{itemize}\setlength\itemsep{0.8em}
  \item Rules prohibiting one-sentence paragraphs and AI-ish
        throat-clearing were added to the style guide.
  \item The textbook uses second person; introductory cases use
        third-person narrative---a deliberate split the style guides enforce.
  \item When the model slipped into pedantic ``landscape of risk'' or
        ``your risk management journey'' language, I deleted it.
  \item When it produced competent but generic examples, I substituted
        firms and episodes I could actually teach about.
\end{itemize}

\bigskip
\textit{I should have spent time reading and learning what engaging writing
actually is. I still have a copy of Strunk and White on my bookshelf. I
probably have not looked at it since college.}

\end{frame}

% -----------------------------------------------------------------------
\section{Implications for Academic Authorship}
% -----------------------------------------------------------------------

\begin{frame}{Speed, Scope, and Accountability}

\textbf{Speed and scope}

\smallskip
The nine-chapter revision cycle compressed work that would otherwise have
spanned many months into about a week for the first draft, a semester
teaching with it, and a week to revise. That is a real productivity gain.

\bigskip
\textbf{Accountability}

\smallskip
Disclosing AI assistance is necessary; treating the model as a co-author
is not. The model has no reputation to lose when the Boeing case dates are
wrong or the present-value calculation is wrong. I remain the author; the
AI is infrastructure.

\bigskip
\textit{When it starts to demand royalties, we can talk.}

\end{frame}

% -----------------------------------------------------------------------
\begin{frame}{Limits of AI Collaboration}

The AI did not reliably:

\begin{itemize}\setlength\itemsep{0.7em}
  \item choose what belongs in a nine-chapter undergraduate treatment;
  \item resist generating plausible nonsense citations;
  \item know Iowa institutional formatting preferences without being told;
  \item finish long integration jobs without monitoring---someone must still
        ask ``are you stuck?'';
  \item stay neutral on topics with a political overtone---it would gush
        about the virtues of ESG. I had to push back to ensure neutrality.
\end{itemize}

\end{frame}

% -----------------------------------------------------------------------
\begin{frame}{A Modest Proposal for Departments}

If business schools desire faculty to produce original, on-point textbooks,
three policies would help:

\begin{enumerate}\setlength\itemsep{0.8em}
  \item Count verified open texts toward teaching/service portfolios when
        they include reusable assets (slides, code, style packages).
  \item Treat AI disclosure in syllabi and prefaces as standard academic
        practice rather than a confession.
  \item Fund \textbf{verification time}---copy editing, fact checking, and
        student piloting.
\end{enumerate}

\bigskip
\textit{The marginal cost of draft prose has collapsed while the marginal
cost of editing and verifying has not.}

\end{frame}

% -----------------------------------------------------------------------
\begin{frame}{Conclusion}

This edition of \textit{Enterprise Risk Management} exists because large
language models made first drafts cheap and a human author made the more
expensive final judgments. This makes sense.

\bigskip
For those considering similar work:

\begin{itemize}\setlength\itemsep{0.6em}
  \item Invest early in style guides and repository structure.
  \item Treat assertions and citations as false until verified.
  \item Box the boring framework sections.
  \item Write the preface yourself.
\end{itemize}

\bigskip
The book is openly available at:\\
\texttt{github.com/ProfMFGrace/erm-book-2ed}\\
(CC BY-NC 4.0)

\end{frame}

% -----------------------------------------------------------------------
\begin{frame}[plain]
  \begin{tikzpicture}[remember picture, overlay]
    \fill[ERMIowaBlack] (current page.south west)
                         rectangle (current page.north east);
    \node[align=center, text=ERMIowaGold,
          font={\usebeamerfont{title}}]
      at (current page.center) {Thank You};
    \draw[ERMIowaGold, line width=1.2pt]
      ([yshift=-0.8cm, xshift= 2.5cm]current page.center) --
      ([yshift=-0.8cm, xshift=-2.5cm]current page.center);
    \node[align=center, text=white,
          font={\usebeamerfont{subtitle}}]
      at ([yshift=-1.6cm]current page.center)
      {\texttt{mfgrace@uiowa.edu}};
  \end{tikzpicture}
\end{frame}

\end{document}
